% sections/intro.tex -- introduction.
\section{Introduction}\label{sec:intro}

The Betti numbers $\beta_0,\beta_1,\beta_2,\beta_3$ of a simplicial complex count
its connected components, independent loops, enclosed voids and, in dimension
three, independent $3$-cycles; over a field they are determined by the ranks of
the boundary matrices~\cite{munkres1984elements,hatcher2002algebraic,edelsbrunner2010computational}.
In many uses the complex is not fixed. Adaptive meshes are refined and
coarsened by local moves~\cite{pachner1991pl}, cavities are carved into a volume
mesh and refilled, holes are opened in a surface and closed again, and the
proximity complex of a point cloud changes as the points
move~\cite{carlsson2009topology}. After each modification one wants to know
whether it changed the topology. Recomputing the Betti numbers from scratch
answers this, but it reads the whole complex after every modification, however
small the modification is.

Over a field, $\beta_k=n_k-r_k-r_{k+1}$, where $n_k$ is the number of
$k$-simplices and $r_k$ the rank of the $k$-th boundary matrix~$\Bm{k}$.
Maintaining Betti numbers therefore means maintaining ranks, and the problem has
more structure than general dynamic rank: the boundary of a simplex does not
change while the simplex exists, so each atomic edit inserts or deletes one
column of one boundary matrix, and rows appear and disappear only while they are
zero (\cref{obs:structural}). Insertions alone are handled by reducing the new
column against the existing pivots: each inserted simplex either creates a cycle
or destroys one~\cite{delfinado1995incremental}. Deletions are the difficulty.
After reduction, a column is a sum of original columns; when one of these is
deleted, every reduced column that absorbed it becomes invalid, and the reduced
matrix alone does not record which columns these are. This paper maintains that
record and studies what it costs.

\paragraph{Contributions.}
\begin{enumerate}[label=(\roman*),leftmargin=2.2em]
\item \emph{Exact dynamic reduced basis maintenance.} For each dimension we
  maintain a reduced matrix $R=\Bm{k}V$ together with the change of basis $V$
  and, for every simplex, the columns of $V$ that contain it. Insertion is the
  standard incremental reduction; a deletion corrects every basis chain that
  contains the deleted simplex, by a cycle (Case~A) or by the column with the
  smallest pivot (Case~B). We prove that the Betti numbers over $\Z$ are exact
  after every edit, for every order of the cells and in every dimension, and
  that initialisation and rebuilds by the static reduction with clearing produce
  a valid state (\cref{thm:exact,cor:rebuild}). The deletion step is the
  rank-only specialisation of the deletion updates used for zigzag persistence
  and warm-started $R=DV$
  factorisations~\cite{carlsson2009zigzag,luo2023warmstart,giunti2026pruning};
  we give a self-contained proof.
\item \emph{A Morse-induced order and its properties.} Ordering the cells by
  reverse removal time in a collapse makes every gradient pair an apparent pair
  (\cref{lem:apparent}, close to known results), so a build performs column
  additions only in the columns of critical cells (\cref{cor:critical-work}), and
  on those columns it performs exactly the reduction of the Morse complex
  (\cref{cor:morse-reduction}). Between rebuilds the order is a heuristic, and
  we measure its effect.
\item \emph{Output-sensitive costs and the limits of locality.} The cost of an
  edit is stated in terms of footprints and reduction chains
  (\cref{prop:cost}). On two explicit families, a cycle graph and a
  triangulated sphere, every toggle of one simplex costs this data structure
  $\Omega(n)$ operations (\cref{prop:nonlocal}), so no bound in terms of the size
  of an edit holds.
\item \emph{An evaluation that identifies where the method is effective.} On
  five workload families (punctured surfaces; refinement, coarsening and
  cavities in a $3$-ball; flag complexes of moving points; non-locality
  witnesses) with up to \NumSizeMax{} simplices, we compare the method with
  static recomputation in insertion order, in a canonical order that cannot
  drift with the history, and in GUDHI~\cite{maria2014gudhi}, and with ablations
  of the order and of the rebuild policy (\cref{sec:experiments}). Every Betti
  vector reported by an exact method was checked against at least one reference
  computation, \NumCoverageExternalPct\,\% against one outside our engine; none
  disagreed. Where edits create and destroy only small cycles, the update
  latency hardly depends on the size of the complex: on the largest refinement
  instance (\NumWTwoNThirtyThreeSEightSize{} simplices) every variant of the
  dynamic reduced basis is
  \NumWTwoNThirtyThreeSEightDrbSpeedupFullIdMin--\NumWTwoNThirtyThreeSEightDrbSpeedupFullIdMax{}
  times faster than static recomputation, with or without the Morse order. On
  punctured surfaces it is at most \NumWOneSphereLFiveMmhmSpeedupFullId{} times
  faster on average, some variants are slower, every variant has a worse tail
  latency, and combinatorial counters are far cheaper; on the non-locality
  witnesses its cost grows linearly with the complex. At build time the Morse
  order saves nothing against the structured insertion orders of our generators,
  although any structured order is far better than a random one; rebuilding pays
  off where the basis grows, as on cavities, and not everywhere.
\end{enumerate}

\paragraph{Scope.}
All statements are over $\Z$; we say nothing about integer homology or torsion.
The maintained basis yields ranks and a basis of each cycle space, but no
persistence pairing. We prove no bound on the cost of an update in terms of its
size, and we make no sparsity claim for the state between rebuilds.

\paragraph{Outline.}
\Cref{sec:overview} states the problem and gives an overview of the algorithm.
\Cref{sec:drb} defines the dynamic reduced basis and proves its exactness,
\cref{sec:morse} introduces the Morse order and its properties, and
\cref{sec:cost-locality} analyses the cost of an update and proves that no
locality bound holds. \Cref{sec:experiments} reports the experiments,
\cref{sec:discussion} discusses limitations and when to use which method,
\cref{sec:related} reviews related work, and \cref{sec:conclusion} concludes.
